%------------------------------------------------
% Latex-Grundger�st f�r Seminarausarbeitungen
%
% zu Erzeugen mit
%   pdflatex ausarbeitung.tex
%
% von Carsten Gutwenger
% angepasst von Igor Vatolkin
%------------------------------------------------

\documentclass[a4paper,12pt,twoside]{article}
\usepackage[a4paper,left=3.5cm,right=2.5cm,bottom=3.5cm,top=3cm]{geometry}

% Encoding der .tex-Datei
% muss je nach Editor-Einstellungen eventuell angepasst werden!
%   latin1    Unix/Windows
%   ansinew   Alternative f�r Windows
%   applemac  Apple
%   utf8      Alternative f�r Unix/Windows/Apple
\usepackage[latin1]{inputenc}

\usepackage[ngerman]{babel}                                 % deutsche Sprache
\usepackage[T1]{fontenc}                                    % Unterst�tzung f�r Umlaute mit Fonts
\usepackage[pdftex]{graphicx,color}                         % Einbetten von Grafiken, Farbe
\usepackage[format=plain,small,bf,indention=7.5mm]{caption} % verbesserte Beschriftung von Abbildungen
\usepackage{subfig}                                         % subfigures
\usepackage{url}                                            % \url{}-Kommando
\usepackage{fancyhdr,float}                                 % Kopf-/Fu�zeilen
\usepackage[pdftex,pdfpagelabels]{hyperref}                 % Hyperlinks im PDF
\usepackage{natbib,doi}                                     % Literaturverzeichnis, DOIs
\usepackage[section,boxed]{algorithm}                       % Floating-Umgebung f�r Algorithmen
\usepackage{algpseudocode}                                  % Pseudocode f�r Algorithmen

\captionsetup{farskip=10pt,topadjust=0pt,captionskip=10pt,nearskip=0pt,margin=10pt}

\pagestyle{fancy}
\fancyhead[LE,RO]{\thepage}
\fancyhead[RE]{\nouppercase{\slshape \leftmark}}
\fancyhead[LO]{\nouppercase{\slshape \rightmark}}
\fancyfoot[C]{}


\begin{document}

%------------------------------------------------
% Titelseite
%------------------------------------------------

\begin{titlepage}
\vspace*{-2cm}
\newlength{\links}
\setlength{\links}{-1.5cm} \sf \LARGE

\hspace*{\links}
\begin{minipage}{12.5cm}
\includegraphics[width=8cm]{tud_logo_rgb}
\end{minipage}

\vspace*{4cm}

\large
\begin{center}
{\Large Seminarausarbeitung} \\[1ex]
{\LARGE\textbf{Darwin's Avatars}}\\[3ex]
Marian Kleineberg, Christopher Riesner\\[1ex]
\today\\[7ex]
im Rahmen des Proseminars\\[1ex]
{\Large\textbf{Computational Intelligence in Spielen}}\\[1ex]
von Dr.~Igor Vatolkin\\[1ex]
Wintersemester 2016/17
\end{center}

\vspace*{5cm}
\hspace*{\links}
\begin{minipage}[b]{15cm}
\normalsize \raggedright

\textbf{Basierend auf:}\\
D.Lessin, S.Risi, Darwin's Avatars: a Novel Combination of Gameplay and Procedural Content Generation, GECCO '15, July 11 - 16
\end{minipage}

\definecolor{TUGreen}{rgb}{0.517,0.721,0.094}
\vfill
\hspace*{\links}
\begin{minipage}[b]{8cm}
\normalsize \raggedright
Fakult�t f�r Informatik\\
Lehrstuhl f�r Algorithm Engineering (Ls11)\\
Technische Universit�t Dortmund\\
\url{http://ls11-www.cs.tu-dortmund.de}
\end{minipage}

\end{titlepage}

%------------------------------------------------
% Inhaltsverzeichnis
%------------------------------------------------

\tableofcontents
\clearpage


%------------------------------------------------
% Einf�hrung
%------------------------------------------------

\section{Einf�hrung}

In dieser Ausarbeitung wird ein Spiel betrachtet, in dem es darum geht f�r den Spieler fremde Kreaturen 
zu steuern. Damit der Spieler stets f�r ihn unbekannte Kreaturen kontrollieren kann wird eine Methode ben�tigt, mit der neue Kreaturen erzeugt werden k�nnen. Dies wird erm�glicht durch die prozedurale Generierung solcher Kreaturen. Dazu wird ein evolution�rer Algorithmus verwendet. Durch die Vielfalt der daraus resultierenden Kreaturen entsteht ein langfristiger Spielspa�.

Dar�ber hinaus wird im evolution�ren Algorithmus ein neuronales Netz entwickelt, welches einer KI die
Steuerung der Kreatur erm�glicht. Es ergibt sich die M�glichkeit gegen eine KI, oder gegen einen weiteren menschlichen Spieler anzutreten.

Der evolution�re Algorithmus entstammt einer fr�heren Arbeit von Lessin \cite{LeFuMi13}. Er berechnet f�r jede Kreatur einen Fitnesswert, der ihre Bewegungsgeschwindigkeit und �sthetik bewertet. �ber mehrere hundert Generationen werden die Kreaturen mit dem h�chsten Fitnesswert selektiert und durch Paarung oder Mutation neue Kreaturen generiert.

%------------------------------------------------
%Beispiel f�r Referenzen: \cite{HanKno06} (Zeitschriftenartikel) und \cite{TarekEtAl07} (Konferenzartikel).
%------------------------------------------------

%------------------------------------------------
% Literaturverzeichnis
%------------------------------------------------

\bibliographystyle{plain}
\bibliography{literatur}
\addcontentsline{toc}{section}{\bibname}


\end{document}
